\subsection{Solucion desarrollada}

\begin{respuesta}{Capacidad del proceso}
Sea \(F\) el flujo de arandanos frescos y \(M\) el flujo de masa base. Las restricciones son:
\[
F\le 300,\qquad 0.8F\le 280,\qquad 0.2F\le 65,\qquad M\le 600.
\]
La mezcla recibe:
\[
M+0.98F\le 800,
\qquad
M+0.98F\le 850.
\]
Como la masa base puede llegar a \(600\) kg/h, usamos \(M=600\). La mezcladora queda activa:
\[
600+0.98F=800
\Rightarrow F=204.08\ \text{kg/h}.
\]
La salida final despues del control de calidad es:
\[
0.97(800)=776\ \text{kg/h de galletas}.
\]
En un ano:
\[
\text{arandanos frescos}=204.08(8)(240)=391836.73\ \text{kg/ano},
\]
\[
\text{masa base}=600(8)(240)=1152000\ \text{kg/ano}.
\]

Si se duplican dos maquinas, conviene duplicar mezclado y cocido. Con eso:
\[
F=300,\qquad M=600,\qquad M+0.98F=894.
\]
La capacidad final queda:
\[
0.97(894)=867.18\ \text{kg/h}.
\]
\end{respuesta}

\begin{respuesta}{Envases: periodo fijo y revision continua}
Tomamos \(z=1.645\), \(\sigma=\sqrt{2500}=50\), \(D=50000\) unidades/mes y \(L=0.25\) meses.

Para revision periodica mensual:
\[
S_T=50000(1+0.25)+1.645(50)\sqrt{1.25}=62591.96.
\]
La politica es ordenar mensualmente hasta una posicion de inventario de \(62592\) envases, de modo que el pedido llegue antes del mes cubierto.

Para revision continua:
\[
Q^*=\sqrt{\frac{2(50000)(90000)}{1.5}}=77459.67,
\]
\[
ROP=50000(0.25)+1.645(50)\sqrt{0.25}=12541.13.
\]
Se ordenan aproximadamente \(77460\) envases cada vez que la posicion baja a \(12541\).
\end{respuesta}

\begin{respuesta}{Wagner-Whitin}
Demandas:
\[
D_1=49000,\qquad D_2=53000,\qquad D_3=47000.
\]
Con \(S=90000\) y \(h=1.5\):
\[
C(1,1)=90000,\quad C(1,2)=90000+1.5(53000)=169500,
\]
\[
C(1,3)=90000+1.5(53000)+2(1.5)(47000)=310500.
\]
La programacion dinamica entrega:
\[
F_1=90000,\qquad F_2=169500,\qquad F_3=250500.
\]
La decision optima es:
\[
Q_1=49000+53000=102000,\qquad Q_2=0,\qquad Q_3=47000.
\]
Es decir, ordenar en el mes 1 para cubrir meses 1 y 2, y volver a ordenar en el mes 3.
\end{respuesta}
